\subsection{行列式的计算}

引理1. 设 A 是一个交换环，M 是一个以 \((e_j)_{j \in J}\) 为基的自由 A-模，其中指标集 \(J\) 是全序的。对于每个整数 \(p \leqslant \operatorname{Card}(J)\) 、每个交错 \(p\) -线性函数 \(f: \mathbf{M}^p \to \mathbf{N}\) （其中 \(N\) 是一个 A-模）以及每一族 \(p\) 个元素 \(x_i = \sum_{j \in J} \xi_{ji} e_j\) （属于 \(M\) ； \(1 \leqslant i \leqslant p\) ），

\[
\begin{array}{r l} f (x _ {1}, x _ {2}, \dots , x _ {p}) & = \sum_ {j _ {1} <   j _ {2} <   \dots <   j _ {p}} \left(\sum_ {\sigma \in \mathfrak {S} _ {p}} \varepsilon_ {\sigma}. \xi_ {j _ {\sigma (1)}, 1} \xi_ {j _ {\sigma (2)}, 2} \dots \xi_ {j _ {\sigma (p)}, p}\right) f (e _ {j _ {1}}, \dots , e _ {j _ {p}}) \end{array}\tag{10}
\]

其中 \((j_k)_{1\leqslant k\leqslant p}\) 遍历由 J 的元素组成的 \(p\) 项严格递增序列的集合。

现在

\[
f (x _ {1}, \dots , x _ {p}) = \sum_ {(j _ {k})} \xi_ {j _ {1}, 1} \xi_ {j _ {2}, 2} \dots \xi_ {j _ {p}, p} f (e _ {j _ {1}}, e _ {j _ {2}}, \dots , e _ {j _ {p}})
\]

其中 \((j_k)_{1 \leqslant k \leqslant p}\) 遍历所有 \(p\) 项序列（ \(J\) 的元素组成的）；然后只需将 §7 第 4 号命题 7 的推论 1 应用于 \(f\) 。

特别地，若J是有限的且含有n个元素，且 \(x_{i} = \sum_{j \in J} \xi_{ji} e_{j} (1 \leqslant i \leqslant n)\) 是 M 的 n 个元素，则

\[
\begin{array}{l} \text {(11)} \quad x _ {1} \wedge x _ {2} \wedge \dots \wedge x _ {n} \\ = \left(\sum_ {\sigma \in \mathfrak {S} _ {n}} \varepsilon_ {\sigma} \xi_ {j _ {\sigma (1)}, 1} \xi_ {j _ {\sigma (2)}, 2} \dots \xi_ {j _ {\sigma (n)}, n}\right) e _ {j _ {1}} \wedge e _ {j _ {2}} \wedge \dots \wedge e _ {j _ {n}} \end{array}
\]

其中 \((j_k)_{1\leqslant k\leqslant n}\) 是 J 的 \(n\) 个元素按递增顺序排列成的唯一序列，从而

\[
\det (x _ {1}, x _ {2}, \dots , x _ {n}) = \sum_ {\sigma \in \mathfrak {S} _ {n}} \varepsilon_ {\sigma}. \xi_ {j _ {\sigma (1)}, 1} \xi_ {j _ {\sigma (2)}, 2} \dots \xi_ {j _ {\sigma (n)}, n}.\tag{12}
\]

采用引理1的记号，比较公式(10)和(12)可得

\[
x _ {1} \wedge x _ {2} \wedge \dots \wedge x _ {p} = \sum_ {\mathrm{H} \in \mathfrak {F} _ {p} (J)} \det (x _ {\mathrm{H}, 1}, x _ {\mathrm{H}, 2}, \dots , x _ {\mathrm{H}, p}) e _ {\mathrm{H}}\tag{13}
\]

其中 \(\mathfrak{F}_{p}(\mathrm{J})\) 是 J 的具有 p 个元素的子集组成的集合，并且对每个子集 \(\mathrm{H}\in\mathfrak{F}_{p}(\mathrm{J})\) ，我们记 \(x_{H,i}=\sum_{j\in H}\xi_{ji}e_{j}\) 和 \(e_{H}=e_{j_{1}}\wedge e_{j_{2}}\wedge\cdots\wedge e_{j_{p}},(j_{k})_{1\leqslant k\leqslant p}\) 是 H 的元素按递增顺序排列成的序列；这里约定 \(\det(x_{\mathrm{H},1},\ldots,x_{\mathrm{H},p})\) 是相对于基 \((e_{jk})_{1\leqslant k\leqslant p}\) 取的。

命题7. 设I为有限集，且 \(X = (\xi_{ji})_{(j, i) \in I \times I}\) 为交换环A上的(I, I)型方阵。则

\[
\det X = \sum_ {\sigma \in \mathfrak {S} _ {I}} \varepsilon_ {\sigma} \left(\prod_ {i \in I} \xi_ {\sigma (i), i}\right)\tag{14}
\]

其中 \(\sigma\) 遍历群 \(\mathfrak{S}_{\mathbf{I}}\) （ \(\mathbf{I}\) 的置换组成的群），而 \(\varepsilon_{\sigma}\) 是 \(\sigma\) 的符号（I，§5，第 7 号）。

只需考虑 \(I = [1, n] \subset N\) 的情形，然后只需应用公式 (12)，其中 \((e_{i})_{1 \leqslant i \leqslant n}\) 是 \(A^{n}\) 的典范基，诸 \(x_{i}\) 是 X 的列（参见第 3 号，公式 (6)）。

特别地，对于三阶矩阵的行列式

\[
X = \left( \begin{array}{c c c} \xi_ {1 1} & \xi_ {1 2} & \xi_ {1 3} \\ \xi_ {2 1} & \xi_ {2 2} & \xi_ {2 3} \\ \xi_ {3 1} & \xi_ {3 2} & \xi_ {3 3} \end{array} \right)
\]

我们有

\[
\begin{array}{r l} \det (X) = & \xi_ {1 1} \xi_ {2 2} \xi_ {3 3} + \xi_ {1 2} \xi_ {2 3} \xi_ {3 1} + \xi_ {2 1} \xi_ {3 2} \xi_ {1 3} \\ & - \xi_ {1 3} \xi_ {2 2} \xi_ {3 1} - \xi_ {1 2} \xi_ {2 1} \xi_ {3 3} - \xi_ {1 1} \xi_ {2 3} \xi_ {3 2}. \end{array}
\]

命题 8. 对于交换环上的每个方阵 X，转置矩阵 \(^{t}X\) 的行列式等于 X 的行列式。

假设 X 是 (I, I) 型的。对于每一对有序的置换 \(\sigma\) , \(\tau\) （ \(S_{I}\) 中的；因为乘法是可交换的），

\[
\prod_ {i \in I} \xi_ {\sigma (i), i} = \prod_ {j \in I} \xi_ {\sigma (\tau (j)), \tau (j)}.
\]

特别地取 \(\tau = \sigma^{-1}\) ；利用 \(\varepsilon_{\sigma^{-1}} = \varepsilon_{\sigma}\) 这一事实，可以看出

\[
\det X = \sum_ {\sigma \in \mathfrak {S} _ {I}} \varepsilon_ {\sigma}. \left(\prod_ {i \in I} \xi_ {i, \sigma (i)}\right),\tag{15}
\]

这证明了该命题。

推论 1. 对于 \(n\) 个向量 \(x_1, \ldots, x_n\) （ \(A^n\) 的），设 \(Y(x_1, \ldots, x_n)\) 表示 \(n\) 阶方阵，其第 \(i\) 行为 \(x_i\) （ \(1 \leqslant i \leqslant n\) ）。则映射

\[
(x _ {1}, \dots , x _ {n}) \mapsto \det (Y (x _ {1}, \dots , x _ {n}))
\]

从 \((\mathbf{A}^n)^n\) 到 \(\mathbf{A}\) 是交错的且 \(n\) -线性的。

推论 2. 对于方阵 \(X\) （交换环 \(A\) 上的、阶为有限的），下列条件等价：

\begin{enumerate}
\def\labelenumi{(\roman{enumi})}
\item
  \(X\) 的行是线性无关的；
\item
  \(X\) 的列是线性无关的；
\item
  det X 在 A 中不是零因子。
\end{enumerate}

这直接由第 3 号命题 6 和上述命题 8 得出。

推论3. 设 \(u\) 是有限维自由 A-模 \(M\) 的一个自同态， \({}^t u\) 是对偶模 \(M^*\) 的转置自同态（II，§2，第 5 号，定义 5）；则

\[
\det (t u) = \det (u).\tag{16}
\]

若 X 是 u 关于 M 的一组基的矩阵，则 \({}^{t}X\) 是 \({}^{t}u\) 关于对偶基的矩阵（II，§10，第 4 号，命题 3）；由于 \(\det(u) = \det(X)\) 且 \(\det(^{t}u) = \det(^{t}X)\) ，结论由命题 8 得出。

